\medskip

\textbf{\large Partie A}

\medskip

Le plan est muni d'un repère orthogonal.

On considère une fonction $f$ définie et dérivable sur $\mathbb{R}$. On note $f'$ sa fonction dérivée. On donne ci-dessous la courbe représentative de la fonction dérivée $f'$.

\begin{center}
\psset{unit=1cm,arrowsize=2pt 3}
\begin{pspicture*}(-7.1,-8)(3.1,3)
\psaxes[linewidth=1.25pt,labelFontSize=\scriptstyle]{->}(0,0)(-7,-7.95)(3,3)
\psgrid[gridlabels=0pt,subgriddiv=1,gridwidth=0.3pt]
\psplot[plotpoints=2000,linewidth=1.25pt,linecolor=red]{-7}{3}{x dup mul 3 x mul sub 1 add 2.71828 x exp mul}
\end{pspicture*}

\end{center}

Dans cette partie, les résultats seront obtenus par lecture graphique de la courbe représentative de la fonction dérivée $f'$. Aucune justification n'est demandée.

\medskip

\begin{enumerate}
\item Donner le sens de variation de la fonction $f$ sur $\mathbb{R}$. On utilisera des valeurs approchées si besoin.
\item Donner les intervalles sur lesquels la fonction $f$ semble être convexe.
\end{enumerate}

\medskip

\textbf{\large Partie B}

\medskip

On admet que la fonction $f$ de la partie $\mathbf{A}$ est définie sur $\mathbb{R}$ par 
\[f(x)=\left(x^{2}- 5 x + 6\right) \mathrm{e}^{x}.\]

On note $\mathcal{C}$ la courbe représentative de la fonction $f$ dans un repère.

\medskip

\begin{enumerate}
\item 
	\begin{enumerate}
		\item Déterminer la limite de la fonction $f$ en $+\infty$.
		\item Déterminer la limite de la fonction $f$ en $-\infty$.
	\end{enumerate}
\item Montrer que, pour tout réel $x$, on a $f'(x)=\left(x^{2}-3 x+1\right) \mathrm{e}^{x}$.
\item En déduire le sens de variation de la fonction $f$.
\item Déterminer l'équation réduite de la tangente $(\mathcal{T})$ à la courbe $\mathcal{C}$ au point d'abscisse 0.
\end{enumerate}

On admet que la fonction $f$ est deux fois dérivable sur $\mathbb{R}$. On note $f''$ la fonction dérivée seconde de la fonction $f$. On admet que, pour tout réel $x$, on a $f''(x) = (x + 1)(x - 2) \mathrm{e}^{x}$.

\begin{enumerate}[resume]
\item 
	\begin{enumerate}
		\item Étudier la convexité de la fonction $f$ sur $\mathbb{R}$.
		\item Montrer que, pour tout $x$ appartenant à l'intervalle $[-1~;~2]$, on a $f(x) \leqslant x + 6$.	
	\end{enumerate}
\end{enumerate}

\bigskip
